\chapter{Введение}

Современное компьютерное зрение развивается не только за счёт увеличения размеров датасетов и вычислительных ресурсов, но и за счёт усложнения архитектур нейронных сетей. За последние годы в этой области сформировались несколько крупных архитектурных направлений. Сначала основную роль играли свёрточные нейронные сети, затем широкое распространение получили глубокие остаточные модели ResNet, а позже в задачах зрения стали активно использоваться визуальные трансформеры и гибридные иерархические архитектуры Swin Transformer. Все эти модели достигали высокой точности на задачах классификации изображений. Но точность не объясняет, что именно происходит внутри модели. Две архитектуры могут давать близкий результат, но приходить к нему через разные внутренние признаки.

В прикладной постановке нейронную сеть часто воспринимают довольно просто: на вход подаётся изображение, на выходе получается класс. Для инженерного использования этого обычно хватает. Для исследовательской задачи — уже нет. Между входом и финальным ответом находится цепочка промежуточных представлений, и каждый уровень этой цепочки меняет исходный сигнал. Ранние слои, согласно тому как устроен CNN, обычно ближе к локальным деталям: границам, текстурам, цветовым переходам. Глубокие уровни работают уже с более собранными признаками — частями объектов, формой, семантическими структурами. Поэтому в данной работе рассматривается не только финальное предсказание, но и то, как меняется само представление изображения по глубине модели.

Одним из способов такого анализа является изучение представлений в частотной области. Изображение можно рассматривать как сигнал, состоящий из компонент разного масштаба. Низкие пространственные частоты соответствуют плавным изменениям яркости, крупным областям и общей форме объекта. Высокие пространственные частоты связаны с резкими переходами, границами, текстурами, мелкими деталями и шумовыми компонентами. Если применить такую точку зрения не только к исходному изображению, но и к промежуточным представлениям нейронной сети, появляется возможность исследовать, какие частотные компоненты сохраняются, подавляются или перераспределяются по мере углубления модели.

Для свёрточных нейронных сетей такая постановка выглядит естественной. Внутренние представления CNN имеют форму двумерных карт признаков, то есть сохраняют пространственную структуру изображения. Кроме того, сама операция свёртки имеет тесную связь с частотной областью: фильтр может усиливать или подавлять различные частотные компоненты сигнала. При этом классические CNN обычно строят иерархию признаков: пространственное разрешение карт постепенно уменьшается, число каналов растёт, а рецептивное поле становится шире. Поэтому можно ожидать, что по мере углубления свёрточной сети высокочастотные локальные детали будут постепенно уступать место более низкочастотным и семантически обобщённым признакам.

Архитектуры ResNet являются одним из наиболее важных представителей CNN-семейства. Их ключевая особенность — остаточные связи, позволяющие эффективно обучать глубокие сети и уменьшать проблему деградации качества при увеличении числа слоёв \cite{he2016deep}. Благодаря этому ResNet удобно использовать для анализа влияния глубины на внутренние представления: в одном семействе можно сравнивать модели разного масштаба, такие как ResNet18, ResNet34, ResNet50 и ResNet101. В данной работе ResNet рассматривается как базовый пример иерархической свёрточной архитектуры.

Визуальные трансформеры устроены принципиально иначе. Vision Transformer разбивает изображение на непересекающиеся патчи, преобразует каждый патч в токен и далее обрабатывает последовательность токенов с помощью блоков самовнимания \cite{dosovitskiy2021image}. Эта идея основана на архитектуре Transformer, изначально предложенной для обработки последовательностей \cite{vaswani2017attention}. В отличие от CNN, ViT не использует явную свёрточную локальность и не строит пространственную иерархию тем же способом. Из-за этого возникает важный вопрос: будет ли ViT демонстрировать ту же частотную динамику, что и CNN, или механизм внимания приводит к иному перераспределению спектральной энергии?

Этот вопрос становится особенно важным на фоне исследований, показывающих, что визуальные трансформеры отличаются от CNN не только архитектурно, но и по характеру внутренних представлений. Например, в работе \cite{raghu2021do} показано, что ViT и CNN по-разному агрегируют пространственную информацию и формируют признаки по глубине. Исследование \cite{park2022how} рассматривает механизм самовнимания как оператор, обладающий низкочастотными свойствами, однако из этого не следует автоматически, что все внутренние представления ViT должны монотонно сжиматься к низким частотам. На итоговый спектр влияет не только attention, но и остаточные связи, нормализации, MLP-блоки, позиционные эмбеддинги и способ организации токенов.

Swin Transformer занимает в этом сравнении отдельное место. Его нельзя отнести ни к классическим свёрточным сетям, ни к обычному ViT без оговорок. От трансформеров Swin наследует токенное представление и механизм внимания, а от CNN — многоуровневую обработку с уменьшением пространственного разрешения. В архитектуре это реализуется через оконное внимание, сдвиг окон между соседними блоками и объединение патчей \cite{liu2021swin}. Поэтому Swin удобен как промежуточный объект анализа: на нём можно отдельно посмотреть, как локальность, иерархия и attention-механизм связаны с частотным составом признаков. Сравнение ResNet, ViT и Swin в этой работе фактически сводится к сравнению трёх способов обработки изображения: свёрточной пирамиды, глобальной токенной модели и локально-иерархического трансформера.

Мотивация работы связана с простой проблемой: accuracy показывает качество ответа, но почти ничего не говорит о том, за счёт каких признаков этот ответ был получен. Модели с близкой точностью могут по-разному реагировать на шум, изменение текстуры, доменный сдвиг или потерю мелких деталей. Частотный анализ даёт здесь дополнительный язык описания. Через него можно увидеть, сохраняет ли представление вклад локальных изменений или быстро переходит к более сглаженной структуре. Для задач медицинской визуализации, спутниковых снимков, дефектоскопии, текстур и мелких объектов это не абстрактный вопрос: слишком раннее подавление высоких частот может убрать признаки, которые несут полезную информацию. В более прикладном смысле такие наблюдения могут пригодиться при регуляризации, дистилляции, сжатии моделей и анализе устойчивости к возмущениям.

Теоретическая опора исследования — работы о spectral bias, то есть о частотном смещении нейронных сетей. В \cite{rahaman2019spectral} показано, что сети обычно быстрее осваивают низкочастотные компоненты функций, а более резкие, высокочастотные зависимости требуют большего усилия при обучении. Позднее эта линия была развита в исследованиях frequency principle и формальных способов измерения спектрального смещения \cite{xu2025overview,kiessling2022computable}. Для компьютерного зрения есть и отдельная мотивация: высокочастотные компоненты изображений связывают как с обобщающей способностью, так и с робастностью моделей \cite{yin2019fourier,wang2020highfrequency}. В этой работе такие результаты используются не как прямой предмет проверки, а как основание для анализа: если частоты влияют на поведение модели, то спектр внутренних представлений может многое сказать о том, какие признаки сохраняются по мере прохождения изображения через сеть.

При этом в существующей литературе разные стороны этой проблемы чаще изучаются раздельно. Одни работы рассматривают общие свойства spectral bias и роль частотных компонент в обучении нейронных сетей. Другие анализируют связь частот с робастностью и обобщением CNN \cite{yin2019fourier,wang2020highfrequency}. Отдельное направление посвящено визуальным трансформерам: их механизму внимания, устойчивости и использованию высокочастотной информации \cite{park2022how,bai2022improving,naseer2021intriguing,paul2022vision}. Однако для сравнения CNN, ViT и Swin требуется единая экспериментальная постановка, в которой разные архитектуры анализируются по одним и тем же спектральным характеристикам. Именно эту задачу решает данная работа: она строит общий протокол анализа представлений и позволяет сопоставить частотную динамику трёх архитектурных семейств на единой основе.

Цель данной выпускной квалификационной работы — исследовать спектральные характеристики внутренних представлений CNN и визуальных трансформеров разной архитектуры и глубины, а также определить, как тип архитектуры влияет на распределение спектральной энергии по слоям.

Для достижения поставленной цели были сформулированы следующие задачи:

\begin{enumerate}
\item изучить теоретические основы спектрального анализа нейросетевых представлений и связанные исследования по spectral bias, CNN и визуальным трансформерам;
\item описать архитектурные особенности ResNet, ViT и Swin, влияющие на пространственную и частотную структуру признаков;
\item реализовать единый экспериментальный протокол извлечения промежуточных представлений из моделей разных архитектур;
\item адаптировать спектральный анализ к токенным представлениям ViT и Swin через восстановление двумерной решётки патчей;
\item построить радиальные спектральные кривые для выбранных уровней ResNet, ViT и Swin;
\item рассчитать количественные метрики спектральной динамики: спектральный центроид, долю низкочастотной энергии и скорость спектрального сжатия;
\item сравнить спектральную динамику внутри каждого архитектурного семейства и между семействами ResNet, ViT и Swin;
\item сопоставить спектральные характеристики с точностью классификации и сформулировать выводы о связи между архитектурой, спектральным профилем и качеством модели.
\end{enumerate}

Объектом исследования являются внутренние представления моделей компьютерного зрения. Предметом исследования являются спектральные характеристики этих представлений: распределение энергии по пространственным частотам, изменение частотного профиля по глубине модели и различия между архитектурными семействами.

В качестве экспериментальных объектов рассматриваются три семейства моделей. Семейство ResNet используется как представитель свёрточных архитектур. Семейство Vision Transformer используется как представитель моделей, основанных на глобальной обработке последовательности патчей. Семейство Swin Transformer используется как представитель иерархических трансформеров с локальным оконным вниманием. Эксперименты проводятся на изображениях из набора ImageNet, который является стандартным датасетом для оценки моделей компьютерного зрения \cite{deng2009imagenet}. В практической части используется подвыборка ImageNet 10K \cite{priyerana_imagenet10k}.

В работе были сформулированы четыре исследовательские гипотезы.

Первая гипотеза состоит в том, что в свёрточных сетях семейства ResNet по мере углубления модели происходит устойчивое спектральное сжатие. Иными словами, ожидается снижение спектрального центроида и рост доли низкочастотной энергии от ранних слоёв к глубоким.

Вторая гипотеза состоит в том, что Vision Transformer не должен полностью повторять монотонную низкочастотную динамику CNN. Поскольку ViT работает с последовательностью патчей и использует самовнимание, MLP-блоки, нормализации и остаточные связи, его спектральная динамика может быть немонотонной.

Третья гипотеза состоит в том, что Swin Transformer должен занимать промежуточное положение между CNN и ViT. За счёт иерархической структуры и patch merging он должен приходить к более низкочастотным финальным представлениям, однако на промежуточных блоках возможна более сложная динамика, чем у классической CNN.

Четвёртая гипотеза состоит в том, что спектральные метрики не должны иметь простой линейной связи с точностью классификации. Они отражают характер внутренних представлений, но не являются самостоятельной заменой accuracy и других метрик качества.

Научная новизна работы заключается в применении единого протокола спектрального анализа к архитектурам различной природы: свёрточным сетям, визуальным трансформерам и иерархическим трансформерам. В отличие от исследований, сосредоточенных только на CNN или только на механизме внимания, в данной работе проводится сопоставление трёх архитектурных семейств на одной методологической основе.

Практическая значимость работы состоит в том, что предложенный подход может использоваться как инструмент диагностики моделей компьютерного зрения. Он позволяет оценивать, насколько быстро модель теряет или сохраняет частотное разнообразие признаков. Это может быть полезно при выборе архитектуры для задач, чувствительных к локальным деталям, а также при разработке методов регуляризации, дистилляции и повышения устойчивости моделей.

Структура работы организована следующим образом. В первой главе формулируются актуальность, цель, задачи, объект, предмет и гипотезы исследования. Во второй главе рассматриваются теоретические основы: понятие внутренних представлений, спектральный анализ, частотные свойства изображений, CNN, ResNet, ViT, Swin и связанные исследования. В третьей главе описывается методика эксперимента: извлечение представлений, построение спектральных кривых и расчёт метрик. В четвёртой главе приводятся результаты экспериментов и их интерпретация. В заключении формулируются основные выводы, ограничения работы и направления дальнейших исследований.
